\chapter{干涉}
\problem{3-1}{在杨氏双孔实验中，孔距为 $0.1\,\mathrm{mm}$，孔与屏幕的距离为 $3\,\mathrm m$，对下列三条典型谱线求出干涉条纹的间距：F蓝线（$486.1\,\mathrm{nm}$），D黄线（$589.3\,\mathrm{nm}$），C红线（$656.3\,\mathrm{nm}$）。}
\answerarea
\problem{3-2}{在杨氏双孔实验中，孔距为 $0.45\,\mathrm{mm}$，孔与幕的距离为 $1.2\,\mathrm m$，测得10个亮纹之间的间距为 $1.5\,\mathrm{cm}$，问光源的波长是多少。}
\answerarea
\clearpage
\problem{3-3}{一微波检测器安装在湖滨高出水面 $0.5\,\mathrm m$ 处。当一颗发射 $21\,\mathrm{cm}$ 波长单色微波的射电星体徐徐自地平线升起时，检测器指出一系列信号强度的极大和极小。当第一个极大出现时，射电星体相对地平线的仰角 $\theta$ 为多少？}
\answerarea
\problemfig[0.4]{3-4}{本题图所示装置是昆克（G. Quincke）用来测声波波长的。管口 $T$ 置于单一声源的声源之前，声波分 $A$、$B$ 两股传播到出口 $O$，其中一股 $B$ 的长度像乐队中的长号那样，可以拉出拉进。当 $A$、$B$ 两股等长时，声音近似保持原有的强度。当 $A$ 逐渐拉开到 $d=16.0\,\mathrm{cm}$ 时，在管口 $O$ 处的声音第一次消失，求此声波的波长。}{3_4}
\answerarea
\clearpage
\problem{3-5}{设菲涅耳双镜的夹角为 $20'$，缝光源离双镜交线 $1.0\,\mathrm{cm}$，接收屏幕与光源经双镜所成的两个虚像联线平行，幕与两镜交线的距离为 $210\,\mathrm{cm}$，光波长为 $600.0\,\mathrm{nm}$。问：
\begin{enumerate}\item 干涉条纹的间距为多少？\item 在幕上最多能看到几根干涉条纹？\item 如果光源到两镜交线的距离增大一倍，干涉条纹有什么变化？\item 如果光源与两镜交线的距离保持不变，而在横向有所移动，干涉条纹有什么变化？\item 如果要在幕上出现有一定衬比度的干涉条纹，允许缝光源的最大宽度为多少？\end{enumerate}
}
\answerarea
\clearpage
\problemfig[0.42]{3-6}{点光源置于薄透镜的焦点，薄透镜后放一双棱镜（见本题图）。设双棱镜的顶角为 $3'30''$，折射率为 $1.5$，屏幕与棱镜相距 $5.0\,\mathrm m$，光波长为 $500.0\,\mathrm{nm}$。求幕上条纹的间距。幕上能出现几根干涉条纹？}{3_6}
\answerarea
\problem{3-7}{设劳埃镜的镜长有 $5.0\,\mathrm{cm}$，幕与镜边缘的距离为 $3.0\,\mathrm m$，缝光源离镜面高度为 $0.5\,\mathrm{mm}$，水平距离 $2.0\,\mathrm{cm}$，光波长为 $589.3\,\mathrm{nm}$。求幕上条纹的间距。幕上能出现几根干涉条纹？}
\answerarea
\clearpage
\problem{3-8}{本题图为梅斯林（L. Meslin）干涉装置，将透镜对剖后再沿光轴方向将两半 $L_1$、$L_2$ 错开一定距离。光点 $S$ 位于光轴上，$S'_1$、$S'_2$ 是它的像。
\begin{enumerate}\item 在图上标出相干光束的交叠区；\item 在交叠区中放一屏幕垂直于光轴，幕上干涉条纹的形状是怎样的？\item 设透镜焦距为 $30\,\mathrm{cm}$，$S$ 与 $L_1$ 的距离为 $60\,\mathrm{cm}$，$L_1$ 与 $L_2$ 的距离为 $8.0\,\mathrm{cm}$，光波长为 $500.0\,\mathrm{nm}$。两像间的中点距离透镜 $L_2$ 有多远？在此处放一屏幕，在其上接收到的亮纹间距为多少？\end{enumerate}\fig[0.65]{3_8}}
\answerarea
\clearpage
\problemfig[0.35]{3-9}{本题图所示为一种利用干涉现象测定气体折射率的原理性结构。在 $S_1$ 孔后面放置一长度为 $l$ 的透明容器，当被测气体注入容器而将空气赶出的过程中幕上的干涉条纹就会移动。由移过条纹的根数即可推知气体的折射率。
\begin{enumerate}\item 设待测气体的折射率大于空气折射率，干涉条纹如何移动？\item 设 $l=2.0\,\mathrm{cm}$，条纹移过20根，光波长为 $589.3\,\mathrm{nm}$，空气折射率为 $1.000276$，求待测气体（氯气）的折射率。\end{enumerate}}{3_9}
\answerarea
\problemfig[0.42]{3-10}{瑞利（Rayleigh）干涉仪的结构和使用原理如下（参见本题图）：以钠光灯作为光源置于透镜 $L_1$ 的前焦点，在透镜 $L_2$ 的后焦面上观测干涉条纹的变动。在两个透镜之间安置一对完全相同的玻璃管 $T_1$ 和 $T_2$。实验开始时，$T_2$ 管充以空气，$T_1$ 管抽成真空，此时开始观测干涉条纹。然后逐渐使空气进入 $T_1$ 管，直到它与 $T_2$ 管的气压相同为止。记下这一过程中条纹移动的数目。设光波长为 $589.3\,\mathrm{nm}$，管长 $20\,\mathrm m$，条纹移动98根，求空气的折射率。}{3_10}
\answerarea
\clearpage
\problem{3-11}{用钠光灯作杨氏双缝干涉实验。光源宽度被限制为 $2\,\mathrm{mm}$，带双缝的幕离缝光源 $2.5\,\mathrm m$，为了在幕上获得可见的干涉条纹，双缝间隔不能大于多少？}
\answerarea
\problem{3-12}{一个直径为 $1\,\mathrm{cm}$ 的发光面元，如果用干涉孔径角量度的话，其空间相干性是多少弧度？如果用相干面积量度，则 $1\,\mathrm m$ 远的相干面积为多大？$10\,\mathrm m$ 远的相干面积为多大？}
\answerarea
\clearpage
\problemfig[0.29]{3-13}{把直径为 $D$ 的细丝夹在两块平玻璃的一边，形成尖劈形空气层（见本题图下方）。在钠黄光（$\lambda=589.3\,\mathrm{nm}$）的垂直照射下形成如图上方所示的干涉条纹。试问 $D$ 为多少？}{3_13}
\answerarea
\clearpage
\problemfig[0.38]{3-14}{块规是机械工里用的一种长度标准，它是一钢质长方体，它的两个端面经过精密抛光，达到相互平行。本题图中 $G_1$、$G_2$ 是同规格的两个块规，$G_1$ 的长度是标准的，$G_2$ 是要校准的。校准的方法如下：把 $G_1$ 和 $G_2$ 放在钢质平台面上，使侧面和底部接触。$G_1$、$G_2$ 上面用一块透明平板 $T$ 压住。如果 $G_1$ 和 $G_2$ 的高度（即长度）不等，略有差别，则在 $T$ 和 $G_1$、$G_2$ 之间分别形成尖劈形空气层。它们在单色光照射下产生等厚干涉条纹。
\begin{enumerate}\item 设入射光的波长是 $589.3\,\mathrm{nm}$，$G_1$ 和 $G_2$ 相隔 $5\,\mathrm{cm}$（即图中的 $l$），$T$ 和 $G_1$、$G_2$ 间干涉条纹的间隔都是 $0.5\,\mathrm{mm}$。试求块规 $G_2$ 和 $G_1$ 的高度之差。怎样判断它们谁高谁矮？\item 如果 $T$ 和 $G_1$ 间干涉条纹的间隔是 $0.5\,\mathrm{mm}$，而 $T$ 和 $G_2$ 间的是 $0.3\,\mathrm{mm}$，则说明什么问题？\end{enumerate}}{3_14}
\answerarea
\clearpage
\problem{3-15}{在图3-27所示的干涉膨胀计中，样品与石英环的高度均为 $1\,\mathrm{cm}$，当温度升高 $100^\circ\mathrm C$ 时，视场中的干涉条纹移动20根。求样品的线膨胀系数。设光波长为 $589.3\,\mathrm{nm}$，石英的线膨胀系数为 $0.35\times10^{-6}/{}^\circ\mathrm C$。
\bookfig[0.43]{book_3_27}{图3-27\quad 干涉膨胀计}
\par{\small 图中 $G$ 为标准石英环，$C$ 为待测柱形样品，$\Pi_1$ 为上方石英平板，$\Pi_2$ 为下方支承平板。温度改变时，样品上端与 $\Pi_1$ 之间楔形空气层的厚度改变，可通过干涉条纹移动测量样品与石英环长度的相对变化。}}
\answerarea
\clearpage
\problem{3-16}{本题图a所示为一种测 PN 结的结深 $x_j$ 的方法。在 N 型半导体基质硅片表面经杂质扩散而形成 P 型半导体区。P 区与 N 区的交界面叫 PN 结，PN 结距表面的深度（即 P 区的厚度 $x_j$）叫做结深，在半导体工艺上需要测定结深。测量的方法是先通过磨角、染色，使 P 区和 N 区的分界线稍微地显示出来，然后盖上半反射膜，在它与硅片之间形成尖劈形空气薄膜。用单色光垂直照射时，可以观察到空气薄膜的等厚干涉条纹。数出 P 区空气薄膜的条纹数目 $\Delta k$ 即可求出结深
\[x_j=\Delta k\frac{\lambda}{2}.\]
由于光在金属或半导体表面反射时相位变化比较复杂，用本方法测量结深 $x_j$ 没有考虑此相位突变，因此测量结果不太精确。更精确的测量方法见图b，半反射膜不是盖在 P 区的上面，而是稍微往上翘一点，观察到的干涉条纹如图b下方所示。试说明：（1）干涉条纹为什么会是这样的？（2）若用 $\lambda=550.0\,\mathrm{nm}$ 的光测得斜干涉条纹的间隔为 $0.20\,\mathrm{mm}$，交界线上两点 $AB$ 间的距离为 $1.1\,\mathrm{mm}$，结深 $x_j$ 为多少？（3）此法比图a所示的方法精确在哪里？\fig[0.78]{3_16}}
\answerarea
\clearpage
\problem{3-17}{测得牛顿环从中间数第5环和第15环的半径分别为 $0.70\,\mathrm{mm}$ 和 $1.7\,\mathrm{mm}$。求透镜的曲率半径。设光波长为 $0.83\,\mu\mathrm m$。}
\answerarea
\problem{3-18}{肥皂膜的反射光呈现绿色，这时膜的法线和视线的角约为 $35^\circ$，试估算膜的最小厚度。设肥皂水的折射率为 $1.33$，绿光波长为 $500.0\,\mathrm{nm}$。}
\answerarea
\clearpage
\problem{3-19}{在玻璃表面上涂一层折射率为 $1.30$ 的透明薄膜，设玻璃折射率为 $1.5$。
\begin{enumerate}\item 对于波长为 $500.0\,\mathrm{nm}$ 的入射光来说，膜厚应为多少才能使反射光干涉相消？这时强度反射率为多少？\item 对波长为 $400.0\,\mathrm{nm}$ 的紫光和 $700.0\,\mathrm{nm}$ 的红光来说，第（1）问所得的厚度在两束反射相干光之间产生多大的相位差？（不考虑色散。）\end{enumerate}
\editorial{原材料第（2）问交叉引用印作“第（2）问”，此处按题意改为“第（1）问”。}}
\answerarea
\clearpage
\problem{3-20}{砷化镓发光管制成半球形，以增加位于球心的发光区对外输出功率，减少反射损耗（见本题图）。已知砷化镓发射光波长为 $930\,\mathrm{nm}$，折射率为 $3.4$。为了进一步提高输出功率，常在球形表面涂敷一层增透膜。
\begin{enumerate}\item 不加增透膜时，球面的强度反射率有多大？\item 增透膜的折射率和厚度应取多大？\item 如果用氟化镁（折射率为 $1.38$）能否增透？强度反射率有多大？\item 如果用硫化锌（折射率为 $2.35$）能否增透？强度反射率有多大？\end{enumerate}\fig{3_20}}
\answerarea
\clearpage
\problem{3-21}{如图3-33用肉眼直接观察薄膜表面的干涉条纹。设瞳孔直径为 $3\,\mathrm{mm}$，与表面相距 $30\,\mathrm{cm}$，视线与表面法线夹角 $30^\circ$，薄膜折射率为 $1.5$。
\begin{enumerate}\item 分别计算膜厚 $2\,\mathrm{cm}$ 及 $20\,\mu\mathrm m$ 两种情况下，点源 $Q_1$、$Q_2$ 在观察点 $P$ 产生的光程差改变量 $\delta(\Delta L)$。\item 如果为了保证条纹有一定的衬比度，要求上述光程差改变量的数量级不能超过多少？以此来估计一下对膜厚 $h$ 的限制。\end{enumerate}
\bookfig[0.48]{book_3_33}{图3-33\quad 干涉条纹衬比度下降的原因}}
\answerarea
\problem{3-22}{证明迈克耳孙干涉仪中圆形等倾条纹的半径与整数的平方根成正比。}
\answerarea
\clearpage
\problem{3-23}{用钠光（$589.3\,\mathrm{nm}$）观察迈克耳孙干涉条纹，先看到干涉场中有12个亮环，且中心是亮的。移动平面镜 $M_1$ 后，看到中心吞（吐）了10环，而此时干涉场中还剩有5个亮环。试求：（1）$M_1$ 移动的距离；（2）开始时中心亮斑的干涉级；（3）$M_1$ 移动后，从中心向外数第5个亮环的干涉级。}
\answerarea
\problem{3-24}{在迈克耳孙干涉仪中，反射镜移动 $0.33\,\mathrm{mm}$，测得条纹变动192次，求光的波长。}
\answerarea
\clearpage
\problem{3-25}{钠光灯发射的黄线包含两条相近的谱线，平均波长为 $589.3\,\mathrm{nm}$。在钠光下调节迈克耳孙干涉仪，人们发现干涉场的衬比度随镜面移动而周期性地变化。实测的结果由条纹最清晰到最模糊，视场中吞（吐）490圈条纹，求钠双线的两个波长。}
\answerarea
\problem{3-26}{在一次迈克耳孙干涉仪实验中，所用的线纹标准具长度为 $0.39\,\mathrm{mm}$，如用镉灯（$643.847\,\mathrm{nm}$）作光源，实验时所测得的条纹变动数目应是多少？}
\answerarea
\clearpage
\problem{3-27}{用迈克耳孙干涉仪进行精密测长，光源为 $632.8\,\mathrm{nm}$ 的氦氖激光，其谱线宽度为 $10^{-4}\,\mathrm{nm}$，整机接收（光电转换）灵敏度可达 $1/10$ 个条纹。求这台仪器测长精度为多少？一次测长量程为多少？}
\answerarea
\problem{3-28}{迈克耳孙干涉仪中的一臂（反射镜）以速度 $v$ 匀速推移，用透镜接收干涉条纹，将它会聚到光电元件上，把光强变化转换为电信号。
\begin{enumerate}\item 若测得电信号的时间频率为 $\nu$，求入射光的波长 $\lambda$。\item 若入射光波长在 $0.8\,\mu\mathrm m$ 左右，要使电信号频率控制在 $50\,\mathrm{Hz}$，反射镜平移的速度应为多少？\item 按以上速度移动反射镜，钠黄光产生电信号的拍频为多少？（钠黄光双线波长为 $589.0\,\mathrm{nm}$ 和 $589.6\,\mathrm{nm}$。）\end{enumerate}}
\answerarea
\clearpage
\problem{3-29}{有两个波长 $\lambda_1$ 和 $\lambda_2$，在 $600.0\,\mathrm{nm}$ 附近相差 $0.0001\,\mathrm{nm}$，要用法布里--珀罗干涉仪把它们分辨开来，间隔 $h$ 需要多大？设反射率 $R=0.95$。}
\answerarea
\problem{3-30}{如果法布里--珀罗干涉仪两反射面之间的距离为 $1.0\,\mathrm{cm}$，用绿光（$500.0\,\mathrm{nm}$）做实验，干涉图样的中心正好是一亮斑。求第10个亮环的角直径。}
\answerarea
\clearpage
\problem{3-31}{设法--珀腔长 $5\,\mathrm{cm}$，用扩展光源做实验，光波波长为 $0.6\,\mu\mathrm m$，问：
\begin{enumerate}\item 中心干涉级数为多少？\item 在倾角为 $1^\circ$ 附近干涉环的半角宽度为多少？设反射率 $R=0.98$。\item 如果用这个法--珀腔分辨谱线，其色分辨本领有多高？可分辨的最小波长间隔有多少？\item 如果用这个法--珀腔对日光进行选频，透射最强的谱线有几条？每条谱线宽度有多少？\item 由于热胀冷缩，引起腔长的改变量为 $10^{-5}$（相对值），求谱线的漂移量（相对值）。\end{enumerate}}
\answerarea
\clearpage
\problemfig[0.39]{3-32}{利用多光束干涉可以制成一种干涉滤光片。如本题图，在很平的玻璃片上镀一层银，在银面上加一层透明膜，例如冰晶石（$3\mathrm{NaF}\cdot\mathrm{AlF}_3$），其上再镀一层银。设银面的反射率 $R=0.96$，于是两个银面之间就形成一个膜层，产生多光束干涉。透明膜的折射率为 $1.55$，膜厚 $h=4\times10^{-5}\,\mathrm{cm}$，平行光正入射。问：（1）在可见光范围内，透射最强的谱线有几条？（2）每条谱线宽度为多少？}{3_32}
\answerarea
\clearpage
